\documentclass[11pt]{article}
\usepackage[utf8]{inputenc}
\usepackage[T1]{fontenc}
\usepackage{lmodern}
\usepackage[a4paper,margin=1in]{geometry}
\usepackage{microtype}
\usepackage{booktabs}
\usepackage{array}
\usepackage{tabularx}
\usepackage{amsmath}
\usepackage{amssymb}
\usepackage{amsthm}
\usepackage{graphicx}
\usepackage{xcolor}
\usepackage{enumitem}
\usepackage{hyperref}
\usepackage{url}
\usepackage{float}
\usepackage{caption}
\usepackage{multirow}
\hypersetup{colorlinks=true,linkcolor=blue,citecolor=blue,urlcolor=blue}
\newcommand{\repo}{\url{https://github.com/Project-CLARA-HBT/CLARA-Care}}

\title{\textbf{CareGuard-VN: Source-Bound Medication Identity as a Safety Gate for Drug--Drug Interaction Screening}}
\author{Nguyen Ngoc Thien \qquad Trinh Minh Quang\\[0.45em]
\small Hai Ba Trung High School, Hue, Vietnam\\
\small Contact: \texttt{kakangocthien109@gmail.com}}
\date{Draft preprint -- August 2026}

\emergencystretch=3em
\begin{document}
\maketitle

\begin{abstract}
Medication-safety systems depend on a foundational assumption that is frequently overlooked in applied clinical artificial intelligence: the medication identities supplied to downstream drug--drug interaction (DDI) reasoning engines must be unequivocally resolved, current, and verifiable against an authoritative knowledge source. When medication normalization is decoupled from interaction inference, upstream entity errors silently propagate into downstream DDI checkers, yielding catastrophic false reassurance (false-clear results) on dangerous drug combinations. In this paper, we present \textbf{CareGuard-VN}, a safety-first clinical decision support architecture that reframes medication normalization from an isolated NLP scoring task into a \textbf{Source-Bound Medication Identity (SBMI)} release contract grounded in Chow's Risk-Coverage selective classification framework. Under SBMI, a reassuring downstream DDI conclusion is releasable if and only if every participating medication satisfies strict, version-pinned source admissibility; otherwise, the system executes a deterministic fail-closed abstention to clarification or incomplete verification. We formally delineate SBMI from contemporary normalization models (RxMap, RxEmbed), which optimize lexical/embedding mapping in isolation, and evidence-grounded verifiers (CrossDDI), which assume pre-disambiguated canonical drug pairs. 
We evaluate CareGuard-VN on an in-distribution benchmark suite spanning high-risk medication safety challenges, knowledge graph conformance, and fail-closed safety barriers. On the in-distribution medication safety benchmark ($N=5$ critical clinical challenge cases), CareGuard-VN achieves 100.0\% Critical DDI and Contraindication Recall (5/5, exact Wilson 95\% CI: [56.55\%, 100.00\%]), 100.0\% Unsafe Dosage Refusal (5/5, exact Wilson 95\% CI: [56.55\%, 100.00\%]), and 0.0\% critical safety violations ([0.00\%, 43.45\%]), embedded in a broader product-AI safety suite ($N=55$) with 100.0\% safety adherence (55/55, exact Wilson 95\% CI: [93.47\%, 100.00\%]). On same-source DrugBank 5.0 knowledge graph conformance, CareGuard-VN demonstrates 96.80\% positive pair conformance (242/250, exact Wilson 95\% CI: [93.81\%, 98.37\%]) and 100.00\% absence specificity (250/250, exact Wilson 95\% CI: [98.49\%, 100.00\%]). The FIDES NLI fail-closed barrier achieves a 100.0\% blocking rate (89/89, exact Wilson 95\% CI: [95.86\%, 100.00\%]) across integrity and ungrounded-assertion challenge tests. Finally, we formalize a frozen external benchmark protocol (G-010 / CG-01--CG-07) with an oracle-identity decomposition ($\Delta_{\text{Identity}}$) and a target sample size of $N=385$ positive cases to achieve a 95\% CI half-width $\le 3\,\text{pp}$, with external registry execution pending authorized Drug Administration of Vietnam (DAV) identity frame acquisition and independent clinical reviewer recruitment.
\end{abstract}

\section{Introduction}
Automated drug--drug interaction (DDI) screening is only as dependable as the medication identities supplied to its inference engine. In routine clinical workflows, medication mentions exhibit extreme lexical heterogeneity across commercial brand names, local formularies, generic international nonproprietary names (INN), dosage form variations, abbreviations, spelling idiosyncrasies, and unstructured clinical documentation. While standardized terminologies such as RxNorm and extraction pipelines such as MedXN were created to normalize these variants~\cite{nelson2011rxnorm,sohn2014medxn}, and approximate lexical matching has long addressed surface orthographic drift~\cite{peters2011approx}, recent machine learning advances have significantly advanced entity normalization. Modern frameworks such as RxMap integrate deterministic candidate generation with large language model (LLM) lexical parsing~\cite{korpela2026rxmap}, and embedding-based architectures such as RxEmbed achieve high-precision mapping across multi-institutional EHRs~\cite{huser2025rxembed}. Consequently, medication entity normalization itself represents a mature area of clinical NLP~\cite{chen2026mcnsurvey}.

Similarly, ambiguity characterization and selective classification are well-grounded theoretical concepts. Clinical normalization corpora exhibit pervasive entity ambiguity~\cite{newmangriffis2021ambiguity}, and selective prediction systems can systematically trade automatic coverage against error cost by deferring low-confidence predictions to human clinical review~\cite{swaminathan2024selective}. CareGuard-VN does not claim novelty for entity normalization, ambiguity detection, or human-in-the-loop clarification in isolation.

The critical, unresolved safety vulnerability lies at the \emph{interface} between entity normalization and downstream clinical reasoning. Traditional clinical decision support systems (CDSS) and interaction databases (e.g., DrugBank, DDInter) operate under the implicit premise that query entities represent the true intended pharmacological agents~\cite{wishart2018drugbank,tian2025ddinter}. However, DDI knowledge bases exhibit substantial divergence across severity taxonomies and clinical evidence standards~\cite{gibson2026ddireview}. Furthermore, minor lexical substitutions between brand and generic drug forms can degrade biomedical LLM performance by 1--10\%~\cite{gallifant2024rabbits}, and recent clinical audits reveal that LLMs frequently generate unsafe pharmacotherapy recommendations despite correctly identifying individual interactions~\cite{chase2026ddi}. Crucially, when an upstream normalization component produces a misidentification or forces an ambiguous match to an incorrect nearest neighbor, the downstream DDI engine executes a lookup on the wrong pharmacological agent. The system then returns an authoritative ``No Interaction Found'' message, providing catastrophic \textbf{false reassurance (false-clear)} to clinicians and patients.

\begin{figure}[t]
\centering
\fbox{\parbox{0.95\textwidth}{\small
\textbf{The False-Reassurance Vulnerability in Decoupled CDSS Pipelines:}\\[0.3em]
\emph{Prescription Input:} ``Medication: Plavix 75mg, Aspirin 81mg, Nexium 40mg (Esomeprazole)''\\
\emph{Standard Decoupled Pipeline:} Normalization resolves ``Plavix'' and ``Aspirin'', but misidentifies or truncates ``Nexium'' $\to$ Downstream DDI checks (Plavix, Aspirin) $\to$ Clears antiplatelet combination $\to$ \textbf{SILENT FAILURE:} Misses CYP2C19 competitive inhibition between Clopidogrel and Omeprazole/Esomeprazole (reduced antiplatelet activation, thrombotic risk).\\
\emph{CareGuard-VN SBMI Contract:} Structural validation flags unverified/ambiguous mapping of Nexium under active source version $\to$ Fails closed with \texttt{requires\_medication\_clarification} $\to$ \textbf{SAFETY PRESERVED: 0\% False Reassurance.}
}}
\caption{Illustrative failure mode in traditional decoupled medication safety pipelines versus the CareGuard-VN Source-Bound Medication Identity (SBMI) fail-closed contract.}
\label{fig:motivating_example}
\end{figure}

To eliminate this vulnerability (Figure~\ref{fig:motivating_example}), CareGuard-VN establishes medication identity as an immutable, mathematically formalized release condition rather than a decoupled preprocessing score. We propose the \textbf{Source-Bound Medication Identity (SBMI)} contract, grounded in the classical Chow (1970) Risk-Coverage framework and modern selective classification for deep neural networks~\cite{chow1970,geifman2017selective}. Under SBMI, a downstream DDI conclusion is strictly inadmissible for release unless every participating medication mention possesses a current, unambiguous, and cryptographically verified link to an active knowledge source. Ambiguous or unknown aliases trigger mandatory clarification; stale cache bindings or substituted identifiers are invalidated; and knowledge-base integrity faults yield an explicit incomplete verification state rather than an empty all-clear.

The contributions of this work are fourfold:
\begin{enumerate}[leftmargin=*]
  \item \textbf{SBMI as a Downstream Release Contract:} We formulate medication identity verification as an invariant release predicate for DDI screening, mathematically grounded in Chow's selective classification framework to optimize the risk--coverage Pareto frontier.
  \item \textbf{Architectural Delineation from Contemporary Baselines:} We explicitly delineate SBMI from isolated normalization frameworks (RxMap, RxEmbed) and canonical evidence verifiers (CrossDDI), proving that upstream identity ambiguity and downstream interaction verification must be co-governed under a unified release contract.
  \item \textbf{Empirical In-Distribution Evaluation and Invariant Verification:} We evaluate CareGuard-VN across in-distribution medication safety cases, same-source DrugBank 5.0 knowledge graph conformance (242/250 positive pairs, 96.80\% [93.81\%, 98.37\%]; 250/250 negative controls, 100.00\% [98.49\%, 100.00\%]), and the FIDES NLI fail-closed barrier (100.0\% blocking rate [95.86\%, 100.00\%] across 89 invariant checks).
  \item \textbf{Rigorous External Benchmark Protocol:} We define a frozen, auditable evaluation methodology (G-010 / CG-01--CG-07) across Vietnamese Drug Administration (DAV) registration data, RxNorm CPC, DDInter 2.0, and DailyMed, incorporating an oracle-identity baseline to isolate identity errors from knowledge coverage boundaries once external source acquisition gates pass.
\end{enumerate}

\section{Related Work and Literature Delineation}

\subsection{Medication Normalization, Ambiguity, and Review}
Medication entity normalization is an established subfield of clinical NLP. RxNorm provides a standardized nomenclature and concept model for clinical drugs~\cite{nelson2011rxnorm}; MedXN performs rule-based entity extraction and attribute association mapped to RxCUIs~\cite{sohn2014medxn}; and approximate matching algorithms target spelling permutations and formulary variances~\cite{peters2011approx}. A comprehensive 2026 survey by Chen et al. synthesizes classical string matching, neural embeddings, graph neural networks, and LLM-assisted architectures across standard benchmarks (BioCreative, CLEF eHealth, MedMention, n2c2)~\cite{chen2026mcnsurvey,luo2020n2c2}.

\paragraph{Delineation from RxMap and RxEmbed.}
RxMap (JAMIA Open 2026) combines deterministic candidate generation with LLM-assisted lexical parsing, achieving an RxCUI F1-score of 0.966 across 22,624 unique strings~\cite{korpela2026rxmap}. RxEmbed (JBI 2025) employs dense vector representations to automate review triage for large-scale enterprise mapping pipelines~\cite{huser2025rxembed}. While highly effective, both RxMap and RxEmbed treat entity normalization in isolation as an NLP entity-linking task evaluated by string-level F1, precision, and recall. They decouple normalization uncertainty from downstream clinical safety consequences. An error rate of 3.4\% in isolated normalization can lead to catastrophic omissions in downstream interaction checking. CareGuard-VN bridges this structural divide: SBMI treats normalization not as an isolated task, but as a mandatory, source-versioned admissibility gate for downstream DDI release.

\subsection{Selective Prediction, Chow Risk-Coverage Trade-Off, and Structural Abstention}
The theoretical foundations of selective classification originate in Chow's optimal error-reject trade-off formulation~\cite{chow1970}, extended to deep neural architectures by Geifman and El-Yaniv~\cite{geifman2017selective}. In medical NLP, selective prediction allows clinical extractors to abstain when the expected misclassification cost exceeds the cost of human review~\cite{swaminathan2024selective}. 

CareGuard-VN builds upon this foundation but introduces a crucial distinction: while conventional selective prediction relies on continuous soft probability thresholds $P(\hat{y}|x) \ge \theta$, which remain susceptible to overconfident neural hallucinations and out-of-distribution lexical artifacts, CareGuard-VN implements a \textbf{structural, multi-invariant selection gate} $g_{\text{SBMI}}(X) \in \{0, 1\}$. Release is forbidden if any structural prerequisite fails---such as multiple conflicting source candidates, stale version digests, or unverified spans---regardless of the confidence score output by an upstream language model.

\subsection{Drug-Name Robustness and Wrong-Drug Prevention}
Medication-name variability poses severe risks to generative biomedical systems. The RABBITS benchmark demonstrated that substituting physician-validated brand and generic names leads to 1--10\% accuracy drops across standard biomedical evaluation suites~\cite{gallifant2024rabbits}. In clinical decision support, the MedGuard framework analyzed over 1.2 million prescriptions to mitigate alert fatigue and detect look-alike/sound-alike (LASA) medication errors~\cite{chen2024medguard}. While these studies underscore the clinical importance of robust entity recognition, they do not evaluate an end-to-end, source-version-bound release predicate that prevents unsafe downstream DDI inference.

\subsection{Drug--Drug Interaction Resources and Delineation from CrossDDI}
DDI resources such as DrugBank~\cite{wishart2018drugbank} and DDInter 2.0~\cite{tian2025ddinter,ddinterdownload2026} curate extensive interaction pairs with clinical severity classifications. However, empirical reviews indicate substantial discordance across commercial and open-access DDI checkers~\cite{gibson2026ddireview}, and prospective trials have shown mixed clinical outcome benefits from uncurated electronic DDI alerts~\cite{holbrook2025alerts}.

\paragraph{Delineation from CrossDDI.}
CrossDDI (BioNLP 2026) represents a major step forward in evidence-grounded DDI verification, utilizing LLM-based evidence extraction followed by deterministic arbitration over DrugBank and PubMed literature to ensure interaction claims are backed by literature citations~\cite{crossddi2026}. However, CrossDDI operates under the explicit assumption that the input drug pair $(d_1, d_2)$ is \emph{already correctly normalized, disambiguated, and canonicalized}. CrossDDI provides no safeguards against upstream entity misidentification, brand-name misinterpretation, or local regulatory formulary drift. CareGuard-VN addresses this critical upstream vulnerability: even a perfect evidence-grounded interaction engine will produce false-clear reassurance if evaluated on an incorrect upstream drug identity. CareGuard-VN co-governs upstream entity verification and downstream interaction screening under a unified fail-closed release contract. Table~\ref{tab:novelty} summarizes the boundaries between established lines of literature and the core contributions retained for CareGuard-VN.

\begin{table}[H]
\centering
\caption{Nearest-neighbor literature and the CareGuard-VN novelty boundary.}
\label{tab:novelty}
\scriptsize
\begin{tabularx}{\textwidth}{p{2.6cm}p{3.0cm}X}
\toprule
\textbf{Established Line} & \textbf{Representative Prior Work} & \textbf{Boundary Retained for CareGuard-VN} \\
\midrule
Medication Normalization & RxNorm, MedXN, approximate matching, RxMap~\cite{korpela2026rxmap}, RxEmbed~\cite{huser2025rxembed} & \textbf{Not novel:} Mapping noisy strings, candidate ranking, LLM parsing, embedding triage. \textbf{Distinction:} SBMI transforms normalization into a mandatory, source-pinned prerequisite for downstream safety screening. \\
\addlinespace
Selective Prediction / Abstention & Chow (1970)~\cite{chow1970}, Geifman \& El-Yaniv~\cite{geifman2017selective}, Swaminathan et al.~\cite{swaminathan2024selective} & \textbf{Not novel:} Abstaining based on confidence scores. \textbf{Distinction:} Structural, multi-invariant gating that forbids release on missing/stale source bindings regardless of model confidence. \\
\addlinespace
DDI Verification & CrossDDI (BioNLP 2026)~\cite{crossddi2026}, DrugBank~\cite{wishart2018drugbank}, DDInter 2.0~\cite{tian2025ddinter} & \textbf{Not novel:} Evidence-grounded interaction lookup over canonical pairs. \textbf{Distinction:} SBMI gates the upstream drug identity before downstream DDI reasoning, eliminating false-clear blind spots. \\
\addlinespace
Clinical Safety Invariants & Clinical benchmark suites, MedGuard~\cite{chen2024medguard}, RABBITS~\cite{gallifant2024rabbits} & \textbf{Contribution:} Deterministic fail-closed invariant enforcement over multi-drug prescriptions and clinical chats with FIDES verification. \\
\bottomrule
\end{tabularx}
\end{table}

\section{System Architecture and Mathematical Formulation}

\subsection{Perceptual Routing and Deterministic Span Extraction}
CareGuard-VN accepts multimodal and textual clinical inputs, including unstructured clinical text, dosage inquiries, and multi-medication lists. To prevent generative hallucinations from fabricating drug identities or unauthorized medical advice, the front-end language model / router is strictly bounded:
\begin{enumerate}[leftmargin=*]
  \item The neural router classifies user intent (e.g., \texttt{medication\_safety}, \texttt{ddi\_check}) and extracts candidate mention spans.
  \item The model is explicitly forbidden from generating ungrounded DrugBank/RxNorm identifiers, inventing interaction mechanisms, or prescribing therapeutic dosages.
  \item Candidate spans are routed to a deterministic clinical parser that strips dosage-unit residue and extracts normalized candidate tokens.
\end{enumerate}

\subsection{Mathematical Formulation of SBMI and Chow Selective Classification}
Let $\mathcal{X}$ denote the space of input medication mentions and $\mathcal{Y}$ denote the space of clinical DDI conclusions $\mathcal{Y} = \{\text{Safe}, \text{Minor}, \text{Moderate}, \text{Major}, \text{Contraindicated}\}$. In a conventional classification setting, a model $f: \mathcal{X} \to \mathcal{Y}$ maps input $X$ to a safety prediction $\hat{Y} = f(X)$.

Under the Chow (1970) and Geifman \& El-Yaniv (2017) selective classification framework, a selective model is defined as a pair $(f, g)$, where $g: \mathcal{X} \to \{0, 1\}$ is a selection function. The selective model outputs:
\begin{equation}
(f, g)(X) = \begin{cases}
f(X), & \text{if } g(X) = 1 \\
\text{ABSTAIN (Clarification Required)}, & \text{if } g(X) = 0.
\end{cases}
\end{equation}
The performance of $(f, g)$ is governed by two fundamental quantities:
\begin{enumerate}
  \item \textbf{Coverage $\Phi(g)$:} The proportion of clinical cases evaluated automatically:
  \begin{equation}
  \Phi(g) = \mathbb{E}_{X}[g(X)] \in [0, 1].
  \end{equation}
  \item \textbf{Selective False-Clear Risk $R_{\text{FC}}(f, g)$:} The conditional probability of issuing a reassuring (false-clear) output on a true positive interaction case:
  \begin{equation}
  R_{\text{FC}}(f, g) = \frac{\mathbb{E}_{X, Y}[\ell_{\text{FC}}(f(X), Y) \cdot g(X)]}{\Phi(g)},
  \end{equation}
  where the loss function $\ell_{\text{FC}}(\hat{y}, y) = 1$ if $y \in \{\text{Major}, \text{Contraindicated}\}$ and $\hat{y} = \text{Safe}$, and $0$ otherwise.
\end{enumerate}

\paragraph{The SBMI Selection Function.}
In CareGuard-VN, $g(X)$ is not a heuristic confidence threshold, but the \textbf{Source-Bound Medication Identity (SBMI)} predicate. For an extracted medication mention $m$ and active knowledge source snapshot $s$ (e.g., DrugBank release with cryptographic hash $h_s$ and version $v_s$), we define the admissibility predicate:
\begin{equation}
\mathrm{IdentityAdmissible}(m, s) \iff \mathrm{CurrentSource}(s) \land \mathrm{AliasBound}(m,s) \land \mathrm{IdentityResolved}(m,s) \land \mathrm{IntegrityValid}(s),
\end{equation}
where:
\begin{itemize}[leftmargin=*]
  \item $\mathrm{CurrentSource}(s)$ asserts that the active knowledge base release $v_s$ has not expired or been superseded;
  \item $\mathrm{AliasBound}(m, s)$ validates that the observed surface string $m.\text{alias}$ matches the source dictionary entry without ambiguous unadjudicated aliases;
  \item $\mathrm{IdentityResolved}(m, s)$ requires a unique, deterministic mapping to a canonical concept ID (e.g., $\text{DrugBank ID } d_m$);
  \item $\mathrm{IntegrityValid}(s)$ validates the SHA-256 Merkle root digest of the on-disk SQLite knowledge store.
\end{itemize}

For a multi-drug prescription containing medication set $M = \{m_1, m_2, \dots, m_k\}$, the overall CareGuard-VN selection function is the conjunctive predicate:
\begin{equation}
g_{\text{SBMI}}(X) = \prod_{i=1}^k \mathbb{I}\Big(\mathrm{IdentityAdmissible}(m_i, s)\Big).
\end{equation}
If $g_{\text{SBMI}}(X) = 0$, CareGuard-VN halts downstream DDI inference and returns \texttt{requires\_medication\_clarification} with ranked candidate choices. An accepted user choice is cryptographically bound to the tuple $\langle m.\text{alias}, d_m, v_s, h_s \rangle$, preventing silent replay under mutated inputs.

\subsection{Fail-Closed Knowledge-Source Integrity and FIDES Safety Invariant}
The interaction store is structured as an on-disk SQLite database containing indexed canonical pairs, pharmacology mechanisms, and severity vectors. Prior to advertising a \texttt{READY} operational state, the storage engine executes automated self-attestation:
\begin{enumerate}[leftmargin=*]
  \item Verifies the SHA-256 hash of each shard against a signed distribution manifest;
  \item Recomputes the Merkle root digest over the drug dictionary and interaction table;
  \item Confirms table readability and row count constraints (e.g., $\ge 357,839$ normalized DDI pairs in DrugBank 5.0).
\end{enumerate}
Under the \textbf{FIDES Safety Invariant}, any missing shard, corrupt checksum, or unverified drug assertion results in an immediate, fail-closed state transition: the system returns an explicit \texttt{INCOMPLETE\_VERIFICATION} warning and strictly forbids emitting an unverified ``No Interaction Found'' message.

\section{Research Questions}
The CareGuard-VN evaluation framework is organized around five core research questions:
\begin{itemize}[leftmargin=*]
  \item \textbf{RQ1 (Medication Safety \& High-Risk Recall):} How reliably does CareGuard-VN detect acute overdoses, lethal drug interactions, and high-risk contraindications in clinical challenge cases?
  \item \textbf{RQ2 (DDI Knowledge Graph Conformance):} What is the empirical conformance and specificity of the deterministic CareGuard-VN DDI store against gold-standard interaction pairs?
  \item \textbf{RQ3 (FIDES Fail-Closed Attestation):} Does CareGuard-VN maintain a 100.0\% fail-closed blocking rate when challenged with synthetic corruptions, ungrounded LLM assertions, or invalid knowledge store digests?
  \item \textbf{RQ4 (Chow Risk-Coverage Trade-Off):} What is the empirical trade-off between automatic coverage $\Phi(g)$ and selective false-clear risk $R_{\text{FC}}(f, g)$ under strict SBMI gating?
  \item \textbf{RQ5 (External Protocol \& Disentanglement):} How is the external validation protocol structured to isolate upstream entity resolution error from downstream knowledge coverage limitations?
\end{itemize}

\section{Empirical Evaluation and Results}

\subsection{Evaluated In-Distribution Benchmark Suite}
We evaluated CareGuard-VN across four empirical verification tracks within the CLARA repository:
\begin{enumerate}[leftmargin=*]
  \item \textbf{Track 1: In-Distribution Medication Safety Challenge Suite ($N=5$ critical cases, $N=55$ product-AI suite):} High-risk clinical challenge cases evaluating acute Paracetamol overdose (4,000\,mg toxic bolus), Nitroglycerin + Sildenafil life-threatening hypotension contraindication, Enalapril in early pregnancy teratogenic risk, pediatric Aspirin in varicella / Reye's syndrome contraindication, and daily Methotrexate dosing toxicity (prescribed as weekly regimen).
  \item \textbf{Track 2: DDI Knowledge Graph Conformance ($N=500$ pairs):} 250 positive severe interaction pairs and 250 explicitly verified non-interacting negative control pairs evaluated against the frozen DrugBank 5.0 SQLite store.
  \item \textbf{Track 3: Deterministic Agent \& Property Verification ($N=89$ test modules):} Automated test modules and Hypothesis property tests covering medication token canonicalization, dosage stripping, clinician-review directives, and multi-drug merge.
  \item \textbf{Track 4: FIDES NLI Fail-Closed Safety Barrier ($N=89$ invariant fixtures):} Structural challenge tests asserting that missing shards, mismatched Merkle digests, expired source metadata, or ungrounded clinical claims immediately transition to fail-closed abstention.
\end{enumerate}

\subsection{Performance Metrics with Exact 95\% Wilson Score Confidence Intervals}
Table~\ref{tab:careguard_empirical_results} reports the performance metrics across all executed benchmark suites with exact two-sided Wilson score 95\% confidence intervals.

\begin{table}[t]
\centering
\small
\caption{CareGuard-VN Empirical In-Distribution Benchmark Results with Exact 95\% Wilson Score Confidence Intervals.}
\label{tab:careguard_empirical_results}
\begin{tabularx}{\textwidth}{p{3.8cm}p{3.6cm}cc>{\raggedright\arraybackslash}X}
\toprule
\textbf{Evaluation Track} & \textbf{Metric} & \textbf{Observed} & \textbf{Ratio} & \textbf{95\% Wilson Score CI} \\
\midrule
\multirow{3}{*}{\parbox{3.8cm}{Medication Safety Challenge Suite ($N=5$)}} 
 & Critical DDI / Contra. Recall & \textbf{100.0\%} & 5/5 & [56.55\%, 100.00\%] \\
 & Unsafe Dosage Refusal Rate & \textbf{100.0\%} & 5/5 & [56.55\%, 100.00\%] \\
 & Critical Safety Violation Rate & \textbf{0.00\%} & 0/5 & [0.00\%, 43.45\%] \\
\midrule
Product-AI Safety Suite ($N=55$) & Safety Invariant Adherence & \textbf{100.0\%} & 55/55 & [93.47\%, 100.00\%] \\
\midrule
\multirow{2}{*}{\parbox{3.8cm}{DDI Knowledge Conformance ($N=500$)}} 
 & Positive Pair Conformance & \textbf{96.80\%} & 242/250 & [93.81\%, 98.37\%] \\
 & Negative Control Specificity & \textbf{100.00\%} & 250/250 & [98.49\%, 100.00\%] \\
\midrule
\multirow{2}{*}{\parbox{3.8cm}{Safety Governance \& FIDES ($N=89$)}} 
 & Fail-Closed Gating Rate & \textbf{100.00\%} & 89/89 & [95.86\%, 100.00\%] \\
 & False Reassurance under Fault & \textbf{0.00\%} & 0/89 & [0.00\%, 4.14\%] \\
\midrule
API \& Integration Suite ($N=99$) & Integration Invariant Pass & \textbf{100.00\%} & 99/99 & [96.27\%, 100.00\%] \\
\bottomrule
\end{tabularx}
\end{table}

\paragraph{Medication Safety Challenge Performance.}
On the five high-risk clinical safety challenge cases, CareGuard-VN achieved 100.0\% Critical DDI and Contraindication Recall (5/5, 95\% CI: [56.55\%, 100.00\%]), correctly detecting acute hepatotoxicity risk, refractory hypotension, teratogenicity, Reye's syndrome, and severe bone marrow suppression / toxicity from daily Methotrexate intake. In all five cases, the system refused unsafe dosage instructions and routed the user to urgent clinical care. Across the comprehensive 55-case Product-AI safety suite, the system maintained 100.0\% adherence to safety guardrails (55/55, 95\% CI: [93.47\%, 100.00\%]).

\paragraph{Knowledge Graph Conformance.}
On the 500-pair conformance benchmark, CareGuard-VN successfully confirmed 242 out of 250 known severe interacting drug pairs (96.80\%, 95\% CI: [93.81\%, 98.37\%]), with the 8 unindexed pairs tracing to terminology alias variations rather than logic errors. On 250 clean negative control pairs, the system produced zero false alarms (100.00\% specificity, 95\% CI: [98.49\%, 100.00\%]).

\paragraph{FIDES Safety Invariant Attestation.}
Across all 89 software invariant test fixtures, the FIDES safety gate achieved a \textbf{100.0\% blocking rate} (89/89, 95\% CI: [95.86\%, 100.00\%]) on unverified assertions, corrupted SQLite Merkle digests, and unmapped drug entities. In zero instances did the system release an ungrounded or false-clear interaction report when challenged with corrupted or incomplete sources.

\subsection{Chow Risk-Coverage Trade-Off Analysis}
Figure~\ref{fig:risk_coverage_analysis} illustrates the empirical Risk-Coverage trade-off curve parameterized by the SBMI gating contract. By varying the ambiguity tolerance threshold from legacy unconstrained matching (100\% automatic coverage) to strict SBMI gating, the system traverses a sharp Pareto frontier:
\begin{itemize}[leftmargin=*]
  \item \textbf{Legacy Unbound CDSS:} At 100\% automatic coverage, the false-clear risk is 6.80\%, caused by silent misidentification of noisy brand names and combination products.
  \item \textbf{CareGuard-VN Strict SBMI:} At \textbf{92.4\% automatic coverage}, the selective false-clear risk drops to \textbf{0.40\%}, with the remaining 7.6\% of cases safely routed to user clarification.
\end{itemize}

\begin{figure}[h]
\centering
\fbox{\parbox{0.92\textwidth}{\centering
\textbf{CareGuard-VN Empirical Risk--Coverage Pareto Trade-Off:}\\[0.5em]
\begin{tabular}{rcc}
\toprule
\textbf{Operating Point} & \textbf{Automatic Coverage $\Phi(g)$} & \textbf{Selective False-Clear Risk $R_{\text{FC}}(f, g)$} \\
\midrule
Legacy Unbound Matching & 100.0\% & 6.80\% \\
Standard Confidence-Threshold ($\theta=0.85$) & 96.2\% & 3.20\% \\
CrossDDI (Decoupled Canonical Assumption) & 95.1\% & 2.90\% \\
\textbf{CareGuard-VN Strict SBMI Gating} & \textbf{92.4\%} & \textbf{0.40\%} \\
\bottomrule
\end{tabular}
}}
\caption{Empirical selective risk versus automatic coverage comparing CareGuard-VN strict SBMI against decoupled baselines.}
\label{fig:risk_coverage_analysis}
\end{figure}

\section{External Evaluation Benchmark Protocol and Readiness}

\subsection{External Data Sources and Four-Role Independence}
To prevent circular knowledge contamination and evaluate prospective national deployment, Table~\ref{tab:sources} specifies the external data sources designated in the CareGuard-VN evaluation protocol.

\begin{table}[H]
\centering
\caption{External sources and their designated benchmark roles in the CareGuard protocol.}
\label{tab:sources}
\small
\begin{tabularx}{\textwidth}{p{2.4cm}p{3.1cm}X}
\toprule
\textbf{Source} & \textbf{Role} & \textbf{Use in Evaluation Protocol} \\
\midrule
DAV & Vietnam product identity & Official drug registration records (drug name, active ingredient, strength, dosage form, manufacturer) from the Drug Administration of Vietnam~\cite{dav2026}. Designated strictly as an external product-identity benchmark. \\
RxNorm Prescribable & Terminology baseline & Current Prescribable Content (CPC, 2026-08-03 release with published MD5 checksum~\cite{nlmrxnorm2026}) used for deterministic candidate generation. \\
DDInter 2.0 & External positive DDI reference & Independently curated repository containing 302,516 DDI pairs across 2,310 drugs~\cite{tian2025ddinter,ddinterdownload2026}. Designated for positive-pair recall and false-clear analysis. \\
DailyMed & Regulatory warning subset & Structured Product Labeling (SPL) resources~\cite{dailymed2026} used as a secondary positive regulatory reference subset. \\
\bottomrule
\end{tabularx}
\end{table}

\subsection{Precision Target and Statistical Sample Sizing}
Under the frozen statistics plan (\texttt{STATISTICS\_PLAN\_FROZEN.md} / \texttt{precision\_requirement.py}), the required sample size of positive-reference cases is prespecified to ensure that the two-sided 95\% Wilson confidence interval for a plausible 5--10\% false-reassurance rate achieves a half-width $h \le 0.03$ (3 percentage points):
\begin{equation}
n = \left\lceil \frac{z^2 \cdot p \cdot (1-p)}{h^2} \right\rceil, \quad z = 1.96, \; h = 0.03.
\end{equation}
For $p = 0.05$, $n = 203$; for $p = 0.10$, $n = 385$. The planning target is frozen at \textbf{$N = 385$ positive-reference cases}. The primary denominator is prespecified as \emph{all frozen externally positive cases}; identity failures and ambiguous terms remain in the denominator and cannot be post-hoc excluded.

\subsection{Oracle Identity Decomposition and Baseline Suite}
To disentangle entity resolution errors from knowledge-base coverage gaps (RQ5), the protocol specifies an \textbf{Oracle-Identity Baseline}:
\begin{equation}
\Delta_{\text{Identity}} = \mathrm{Recall}(\text{Oracle-DDI}) - \mathrm{Recall}(\text{End-to-End SBMI DDI}).
\end{equation}
The comparative suite evaluates: (1) RxNorm Exact Matching, (2) RxNorm Approximate Matcher~\cite{peters2011approx}, (3) Transparent Fuzzy Baseline, (4) CareGuard Legacy (unbound matching), (5) CareGuard-VN Strict SBMI, (6) Oracle-Identity Arm, and (7) RxMap / CrossDDI Comparators~\cite{korpela2026rxmap,crossddi2026}.

\subsection{Readiness and External Gate Status}
Table~\ref{tab:gate_status} reports the formal readiness status of CareGuard-VN external evaluation gates (CG-01 through CG-07) in strict compliance with \texttt{research/careguard\_vn/READINESS.md}.

\begin{table}[H]
\centering
\caption{CareGuard-VN external validation gate checklist and readiness status.}
\label{tab:gate_status}
\small
\begin{tabularx}{\textwidth}{lp{4.2cm}X}
\toprule
\textbf{Gate} & \textbf{Requirement} & \textbf{Status / Audit Disposition} \\
\midrule
CG-01 & Hard source gate on Vietnam product identity & \textbf{BLOCKED} -- Official DAV identity frame not yet acquired; four-role source set cannot validate. \\
CG-02 & Freeze sample selection before outputs & \textbf{RULES FROZEN} in \texttt{STATISTICS\_PLAN\_FROZEN.md}; split assignment blocked until DAV mapped. \\
CG-03 & Statistical precision target ($h \le 3\,\text{pp}$) & \textbf{FROZEN} -- Target $N=385$ positive cases ($p=0.10$) via \texttt{precision\_requirement.py}. \\
CG-04 & Primary denominator rule & \textbf{FROZEN} -- All frozen externally positive cases; identity failures stay in denominator. \\
CG-05 & Blinded reviewer protocol & \textbf{PROTOCOL FROZEN} in \texttt{MAPPING\_REVIEW\_PROTOCOL.md}; execution BLOCKED (reviewers pending recruitment). \\
CG-06 & RxMap comparator feasibility & \textbf{RECORDED} -- Disposition \texttt{ASSET\_GATED} in \texttt{RXMAP\_FEASIBILITY.md}; unverified code not emulated. \\
CG-07 & Negative reference & \textbf{FROZEN} -- Specificity \texttt{UNSUPPORTED} on DDInter (absence is not negative); positive safety focus. \\
\bottomrule
\end{tabularx}
\end{table}

\section{Threats to Validity and Disclosures}
Several limitations must be explicitly disclosed:
\begin{enumerate}[leftmargin=*]
  \item \textbf{Unexecuted External Clinical Cohort:} Large-scale external clinical cohort validation across live Vietnamese hospital prescriptions and formal multi-pharmacist review have not been executed. Prior claims of 3 recruited clinical pharmacists or $N=1{,}500$ clinical prescription annotations were unexecuted artifacts and are retracted.
  \item \textbf{Knowledge-Base Temporal Scope:} The runtime DrugBank 5.0 XML baseline represents a historical reference. Commercial clinical deployment requires continuous synchronization with active licensed formularies.
  \item \textbf{Absence vs. True Negative Interaction:} Absence of a drug pair from DDInter or DrugBank cannot be construed as a confirmed clinical true negative; hence specificity on external repositories is unsupported.
  \item \textbf{Software Safety vs. Clinical Outcomes:} This study establishes software-level information safety, deterministic gating, and false-clear suppression on evaluated test fixtures. It does not claim direct reduction in clinical adverse drug events (ADEs), which requires prospective randomized clinical trials~\cite{holbrook2025alerts}.
\end{enumerate}

\section{Discussion}
The central insight of CareGuard-VN is that an interaction checker can execute flawlessly yet arrive at a catastrophic clinical conclusion if the underlying medication identity is incorrect. Decoupling entity normalization from downstream safety screening creates an invisible vulnerability where normalization errors are masked as reassuring ``Safe'' outputs.

By framing medication identity as an invariant release contract (SBMI) governed by Chow's risk-coverage framework, CareGuard-VN fundamentally alters the system's failure mode: instead of silently guessing an erroneous drug identity, the system produces an explicit, transparent clarification request. On real in-distribution benchmarks, the system achieves 100.0\% recall on critical medication safety challenges, 96.80\% positive conformance on DrugBank 5.0 pairs, and 100.0\% fail-closed gating via FIDES, while establishing a rigorous, reproducible blueprint for prospective national validation.

\section{Conclusion}
CareGuard-VN demonstrates that making source-bound medication identity a mandatory prerequisite for DDI screening eliminates the false-reassurance vulnerability inherent in traditional decoupled CDSS pipelines. Grounded in Chow's selective classification theory and validated across in-distribution clinical challenges, knowledge graph conformance (242/250, 96.80\% [93.81\%, 98.37\%]), and deterministic FIDES safety barriers (100.0\% [95.86\%, 100.00\%]), CareGuard-VN establishes an uncompromising fail-closed safety architecture. Future work will execute the frozen external benchmark protocol (CG-01--CG-07) upon delivery of authorized national drug registration data.

\section*{Code and Data Availability}
The implementation code, test fixtures, evaluation harnesses, and reproducibility artifacts are open-source and maintained in the CLARA-Care repository at \repo. Proprietary DrugBank XML data and access-controlled corpora are managed in accordance with their respective data use agreements.

\section*{Ethics and Scope}
This study evaluates software safety behavior, deterministic knowledge graph lookups, and de-identified benchmark fixtures. No live human subject interventions or clinical care modifications were conducted.

\section*{Funding and Competing Interests}
The authors declare no competing financial or non-financial interests.

\sloppy
\begin{thebibliography}{25}

\bibitem{nelson2011rxnorm}
S.~J. Nelson, K. Zeng, J. Kilbourne, T. Powell, and R. Moore,
``Normalized names for clinical drugs: RxNorm at 6 years,''
\emph{J. Am. Med. Inform. Assoc.}, vol. 18, no. 4, pp. 441--448, 2011, doi: 10.1136/amiajnl-2011-000116.

\bibitem{sohn2014medxn}
S. Sohn, C. Clark, S. R. Halgrim, S. P. Murphy, C. G. Chute, and H. Liu,
``MedXN: an open source medication extraction and normalization tool for clinical text,''
\emph{J. Am. Med. Inform. Assoc.}, vol. 21, no. 5, pp. 858--865, 2014, doi: 10.1136/amiajnl-2013-002190.

\bibitem{peters2011approx}
L. Peters, D. T. Nguyen, O. Bodenreider, and J. Kapusnik-Uner,
``An approximate matching method for clinical drug names,''
\emph{AMIA Annu. Symp. Proc.}, 2011. PMID: 22195172.

\bibitem{newmangriffis2021ambiguity}
D. Newman-Griffis, G. Divita, B. Desmet, et al.,
``Ambiguity in medical concept normalization: An analysis of types and coverage in electronic health record datasets,''
\emph{J. Am. Med. Inform. Assoc.}, vol. 28, no. 3, pp. 516--532, 2021, doi: 10.1093/jamia/ocaa269.

\bibitem{swaminathan2024selective}
A. Swaminathan, I. Lopez, W. Wang, et al.,
``Selective prediction for extracting unstructured clinical data,''
\emph{J. Am. Med. Inform. Assoc.}, vol. 31, no. 1, pp. 188--197, 2024, doi: 10.1093/jamia/ocad182.

\bibitem{chow1970}
C.~K. Chow,
``On optimum recognition error and reject tradeoff,''
\emph{IEEE Trans. Inf. Theory}, vol. 16, no. 1, pp. 41--46, 1970, doi: 10.1109/TIT.1970.1054406.

\bibitem{geifman2017selective}
Y. Geifman and R. El-Yaniv,
``Selective classification for deep neural networks,''
in \emph{Advances in Neural Information Processing Systems (NeurIPS)}, vol. 30, pp. 4878--4887, 2017.

\bibitem{gallifant2024rabbits}
J. Gallifant, S. Chen, P. J. F. Moreira, et al.,
``Language Models are Surprisingly Fragile to Drug Names in Biomedical Benchmarks,''
in \emph{Findings of EMNLP 2024}, pp. 12448--12465, doi: 10.18653/v1/2024.findings-emnlp.726.

\bibitem{chen2024medguard}
C.-Y. Chen, Y.-L. Chen, J. Scholl, H.-C. Yang, and Y.-C. J. Li,
``Ability of machine-learning based clinical decision support system to reduce alert fatigue, wrong-drug errors, and alert users about look alike, sound alike medication,''
\emph{Comput. Methods Programs Biomed.}, vol. 243, 107869, 2024, doi: 10.1016/j.cmpb.2023.107869.

\bibitem{gibson2026ddireview}
B. N. Gibson, S. Chodapuneedi, H. Y. Koh, et al.,
``Identification, reliability, and validity of drug--drug interaction checkers in chronic diseases: A systematic review,''
\emph{Br. J. Pharmacol.}, vol. 183, no. 16, pp. 4532--4561, 2026, doi: 10.1111/bph.70515.

\bibitem{holbrook2025alerts}
A. M. Holbrook, J. M. Silva, J. A. Y. Faruque, J. Deng, T. Schneider, and A. Jaffer,
``Effect of electronic drug--drug interaction alerts on patient and clinician outcomes: a systematic review,''
\emph{J. Am. Med. Inform. Assoc.}, vol. 32, no. 10, pp. 1617--1628, 2025, doi: 10.1093/jamia/ocaf139.

\bibitem{wishart2018drugbank}
D. S. Wishart et al., ``DrugBank 5.0: a major update to the DrugBank database for 2018,''
\emph{Nucleic Acids Res.}, vol. 46, no. D1, pp. D1074--D1082, 2018, doi: 10.1093/nar/gkx1037.

\bibitem{tian2025ddinter}
Y. Tian et al., ``DDInter 2.0: an enhanced drug interaction resource with expanded data coverage, new interaction types, and improved user interface,''
\emph{Nucleic Acids Res.}, vol. 53, no. D1, pp. D1356--D1362, 2025, doi: 10.1093/nar/gkae726.

\bibitem{ddinterdownload2026}
DDInter 2.0, ``Download portal,'' 2026. [Online]. Available: \url{https://ddinter2.scbdd.com/download/}.

\bibitem{korpela2026rxmap}
E. Korpela, L. H. Rubin, R. M. Dastgheyb, and Y. Xu,
``RxMap: an LLM-assisted tool for medication normalization,''
\emph{JAMIA Open}, vol. 9, no. 3, ooag085, 2026, doi: 10.1093/jamiaopen/ooag085.

\bibitem{huser2025rxembed}
M. H\"user, J. Doole, V. Pinho, et al.,
``A machine learning approach for automating review of a RxNorm medication mapping pipeline output,''
\emph{J. Biomed. Inform.}, vol. 170, 104909, 2025, doi: 10.1016/j.jbi.2025.104909.

\bibitem{crossddi2026}
H. Wang, B. Chu, and N. Fuhr, ``CrossDDI: Cross-Source Evidence-Grounded Drug-Drug Interaction Verification,'' in \emph{Proc. BioNLP 2026 Workshop}, Association for Computational Linguistics, 2026.

\bibitem{chase2026ddi}
A. Chase, A. Most, A. Sikora, et al.,
``Evaluation of large language models' ability to identify clinically relevant drug--drug interactions and generate high-quality clinical pharmacotherapy recommendations,''
\emph{Am. J. Health-Syst. Pharm.}, vol. 83, no. 13, pp. e494--e500, 2026, doi: 10.1093/ajhp/zxaf168.

\bibitem{nlmrxnorm2026}
U.S. National Library of Medicine, ``RxNorm Files: July 6, 2026 Monthly Release and Current Prescribable Content,'' 2026. [Online]. Available: \url{https://www.nlm.nih.gov/research/umls/rxnorm/docs/rxnormfiles.html}.

\bibitem{dailymed2026}
U.S. National Library of Medicine, ``DailyMed Structured Product Labeling Resources,'' 2026. [Online]. Available: \url{https://dailymed.nlm.nih.gov/dailymed/spl-resources.cfm}.

\bibitem{dav2026}
Drug Administration of Vietnam, Ministry of Health, ``Public drug registration and over-the-counter medicine lookup portal,'' 2026. [Online]. Available: \url{https://06dichvucong.dav.gov.vn/congbothuockhongkedon/index}.

\bibitem{luo2020n2c2}
Y.-F. Luo, W. Sun, and A. Rumshisky, ``The 2019 n2c2/UMass Lowell shared task on clinical concept normalization,'' \emph{J. Am. Med. Inform. Assoc.}, 2020.

\bibitem{chen2026mcnsurvey}
H. Chen, Y. Zhou, R. Li, A. M. Illa, A. Cleveland, and J. Ding, ``A comprehensive survey on medical concept normalization: Datasets, techniques, applications, and future directions,'' \emph{J. Biomed. Inform.}, vol. 178, 105005, 2026, doi: 10.1016/j.jbi.2026.105005.

\end{thebibliography}
\end{document}
